\section{Introduction}\label{sec:introduction}

Let $X$ be a smooth connected projective complex variety, and let
$\widetilde X$ denote its simply connected topological universal cover,
equipped with the lifted complex structure.  Shafarevich's
holomorphic-convexity conjecture asks whether $\widetilde X$ is
holomorphically convex.  Its particularly direct form for varieties with
large fundamental group asks whether $\widetilde X$ is Stein when it
contains no positive-dimensional compact complex analytic subvariety.

Recall that $X$ has \emph{large fundamental group} if, for every
positive-dimensional closed irreducible subvariety $Z\subset X$, the image
of $\pi_1(Z^\nu)\to\pi_1(X)$ is infinite; here $Z^\nu$ is the
normalization.  This is equivalent to the absence of positive-dimensional
compact analytic subvarieties in $\widetilde X$.  We prove the following.

\begin{theorem}\label{thm:main}
There exist a smooth connected projective complex fourfold $X$ and a
simple abelian surface $A\subset X$ such that
\begin{enumerate}
\item $X$ has large fundamental group;
\item $\im\bigl(\pi_1(A)\to\pi_1(X)\bigr)\simeq\Z$.
\end{enumerate}
The universal cover $\widetilde X$ contains no positive-dimensional
compact complex analytic subvariety and is not Stein.
\end{theorem}

The obstruction to Steinness is topological and occurs in a closed
complex surface inside $\widetilde X$.  Since $\pi_1(A)\simeq\Z^4$,
the kernel of the homomorphism in Theorem~\ref{thm:main} has rank three.
The corresponding cover of $A$ is therefore diffeomorphic to
$\R\times(S^1)^3$ and has nonzero third homology.  A Stein surface has
the homotopy type of a CW complex of real dimension at most two.
This contradiction rules out Steinness of the lifted surface and hence
of $\widetilde X$.  Simplicity of $A$ ensures that this topological
obstruction introduces no compact complex curve in that covering.

\subsection{Historical context}
Shafarevich's question extends the uniformization viewpoint for curves
to universal covers of projective varieties.  Koll\'ar's account
\cite[Introduction, Section~0.3]{Kollar1995} explains both the
holomorphic-convexity conjecture and its formulation for large
fundamental groups.  These formulations are related by Remmert
reduction: a holomorphically convex complex manifold admits a proper
holomorphic map to a Stein space with connected compact analytic
fibers \cite[Introduction, p.~1546]{EKPR2012}.  If there are no positive-dimensional
compact analytic subvarieties, all those fibers are points and the
reduction is an isomorphism.  Theorem~\ref{thm:main} therefore also
disproves the holomorphic-convexity conjecture.

An important development was the construction of almost-holomorphic reduction maps by
Campana and Koll\'ar \cite{Campana1994,Kollar1993,Kollar1995}.
They describe, through a very general point, the normalized subvarieties
whose fundamental groups have finite image in the ambient group.
Their existence provides an algebraic
and geometric description of the part of the variety that such
subvarieties fill.  Holomorphic convexity asks in addition for a
proper holomorphic reduction of the universal cover to a Stein space.
The large case makes this additional analytic requirement especially
visible: the expected reduction has no positive-dimensional fibers.

Linearity of the fundamental group supports strong affirmative
results.  Katzarkov and Ramachandran proved holomorphic convexity for
projective surfaces whose fundamental group admits a faithful complex
representation with reductive Zariski closure
\cite{KatzarkovRamachandran1998}.  Eyssidieux treated the reductive
linear case in arbitrary dimension \cite{Eyssidieux2004}.
Eyssidieux, Katzarkov, Pantev, and
Ramachandran proved the linear Shafarevich conjecture for smooth
projective varieties \cite{EKPR2012}.  More recently, Bakker,
Brunebarbe, and Tsimerman developed linear Shafarevich theory for
normal algebraic spaces and quasiprojective varieties
\cite{BBT2024}.  The construction here concerns the unrestricted
fundamental group.  Its infinite-order argument uses compatible
local covers of an arrangement, so it does not require a faithful
linear representation of that group.

A closely related line of inquiry appears in Bogomolov and Katzarkov's
constructions of projective surfaces and their potential counterexamples
to the Shafarevich conjecture
\cite[Section~4]{BogomolovKatzarkov1998}.  Their proposed obstruction
uses infinite chains of compact curves and depends on group-infiniteness
statements left conjectural there.
Here the obstruction is a rank-three lattice cover of a simple abelian
surface.  The group-theoretic task is to prove that its cyclic image
survives all global relations; the local path criterion supplies that
infiniteness argument.

Several established tools make the construction possible.  Arithmetic
ball quotients supply arbitrarily large regions with exactly prescribed
finite hyperplane configurations.  Stover and Toledo's virtual
ramified-cover theorem supplies the required sign covers
\cite{StoverToledo,LlosaPy}, and Zarhin's theorem supplies a
genus-two curve with simple Jacobian \cite{Zarhin}.  Standard
hyperbolic geometry \cite{BridsonHaefliger} underlies the local path
criterion.  Hamm--L\^e's quasi-projective Lefschetz theorem
\cite{HammLe}, together with a deformation argument, allows us to
retain a prescribed surface while passing
from an orbifold to a smooth projective ambient variety.  Each use is
stated with its hypotheses at the relevant point below.

\subsection{The central construction}
The principal task is to embed a simple abelian surface in a smooth
projective variety so that its fundamental group has infinite cyclic
image.  Once such an embedding is available, an ample complete
intersection in a product with a sufficiently large abelian variety
gives Theorem~\ref{thm:main}.  The complete intersection contains the
prescribed surface, and its other fibers over the added abelian variety
are finite.  This separates the two reasons that a subvariety has
infinite fundamental-group image: it either moves in the added abelian
variety or lies in the simple surface.  Section~\ref{sec:reduction}
proves this reduction first.

To obtain the cyclic image, we begin with a genus-two curve $C$ whose
Jacobian is simple.  Its presentation as a double cover of $\PP^1$
uses four branch points in one small disk and two outside that disk.
We call the corresponding meridians red and blue.  A complex hyperbolic
arrangement realizes this marked line as an exceptional fiber.  Extra
hyperplanes meeting the red part of the pencil impose relations that
identify all four red meridians.  The sphere relation then identifies
the two blue meridians.  The image of the genus-two surface group is
therefore cyclic: it lies in the cyclic subgroup generated by the product
of the resulting red and blue involutions.

These relations alone give only an upper bound.  The central difficulty
is to prove that the cyclic image is infinite after all global
identifications and fillings.  We use covers defined separately on
large local models of the arrangement.  We model their monodromies
using the infinite dihedral group generated by two involutions $r,b$;
its even words $(rb)^n$ carry the integer translation label $n$.
These local monodromies need not
define a representation of the full projective fundamental group.
Instead, they agree on the endpoint tests needed to read short paths
in overlapping charts.  The graphs of these local covers are uniformly
Gromov hyperbolic.  A local path criterion then prevents a loop detected
by one chart from becoming trivial through global relations.

The edges of these graphs are paths with bounded \emph{projected
diameter}, rather than bounded length.  Consequently every power of a
fixed loop near the distinguished fiber is a single edge.  The local
criterion detects every nonzero power separately and gives the
infinite-order lower bound.  This feature is essential to the argument;
neither a global coloring nor a family of finite quotients is needed
for that lower bound.

Two copies of the pencil give a distinguished fiber $C\times C$.
An order-four symmetry exchanges the factors and rotates their normal
directions.  Resolving its coarse quotient over this fiber produces a
map from a point blowup of
$\Sym^2 C$, itself the blowup of $\Jac(C)$ at a point.  The same local
argument detects the sum of the two translation labels.  We then
realize the resulting surface map in an ordinary smooth projective
ambient variety and remove its point blowups by projective
modifications.  Throughout these operations we preserve the actual
homomorphism from the surface group.

The local path criterion and the ambient blowdown construction are
formulated separately from the example.  The former converts bounded
chart compatibility and hyperbolicity into nontriviality of loops.
The latter realizes a point blowdown of an embedded surface inside a
new smooth projective ambient variety, preserving the ambient
fundamental group and the induced surface homomorphism.  These two
statements isolate the group-theoretic and projective-geometric
mechanisms of the construction.

\subsection{Structure of the proof}
Section~\ref{sec:paths} proves the local path criterion independently
of the geometric application.  Its hypotheses are bounded-range path
readings, compatibility of endpoint equality, and uniform hyperbolicity
of the reading graphs; no local compactness is required.
Section~\ref{sec:geometry} constructs the arithmetic arrangement,
its sign covers, and the distinguished surface map.
Section~\ref{sec:models} defines the local monodromy labels, proves
their buffered compatibility, and establishes hyperbolicity of the
associated graphs.  Section~\ref{sec:image} verifies the path
criterion for the constructed orbifold and proves that the surface
image is infinite cyclic.  Section~\ref{sec:projective} carries this
image into an ordinary smooth projective variety, removes the point
blowups, and completes the proof using Section~\ref{sec:reduction}.

\subsection{Conventions}
All varieties are over $\C$ unless a field is specified, and all
fundamental groups are taken in the usual topology.  For a smooth
complex orbifold, equivalently a smooth Deligne--Mumford stack in the
constructions below, the fundamental group means that of its
classifying space.  An \emph{ordinary} point has trivial stabilizer.
Fundamental-group image statements are understood up to a common
choice of basepoints and connecting paths.  Projective bundles
parametrize lines.  Complex hyperbolic distance is normalized so that
a point of Euclidean radius $\tanh r$ in the unit ball has distance
$r$ from its center.
